% Standalone source-record documentation; no research attempts.
\documentclass[11pt]{article}
\usepackage[margin=25mm]{geometry}
\usepackage{fontspec,unicode-math}
\setmainfont{Latin Modern Roman}
\setmathfont{Latin Modern Math}
\usepackage{amsmath,amssymb,mathtools,xurl,hyperref}
\hypersetup{colorlinks=true,urlcolor=blue,linkcolor=blue}
\setcounter{secnumdepth}{0}
\setlength{\parindent}{0pt}
\setlength{\parskip}{5pt}
\setlength{\emergencystretch}{3em}
\newcommand{\sref}[2]{\hyperlink{source:#1:#2}{[S#2]}}
\newcommand{\cataloguescope}[1]{\paragraph{Source record.}#1}
\newcommand{\floor}[1]{\left\lfloor#1\right\rfloor}
\begin{document}
\section{10 Lectures and 42 Open Problems — Angular Synchronization via Projected Power Method}\label{statement}

\subsection{Definitions and mathematical statement}

\paragraph{Shared notation.}Standard number systems are $\mathbb N,\mathbb Z,\mathbb Q,\mathbb R,\mathbb C$. Specialized notation and index ranges follow the cited source and the definitions below.


Angular synchronization estimates a vector $z\in\mathbb C^n$ with $|z_i|=1$ from noisy relative-phase measurements. In the standard source model $C=zz^*+\sigma W$, $W$ is Hermitian Gaussian noise. The optimization maximizes $x^*Cx$ subject to $|x_i|=1$, with a common global phase unidentifiable. A projected or generalized power iteration applies the source linear update and normalizes each nonzero coordinate to unit modulus; initialization, shifts and zero-coordinate conventions are part of the algorithm. High probability refers to the source noise law as $n$ grows.
\paragraph{Record scope and qualifications.}
The catalogue title is retained as metadata. Named certificates, contracts, bounds or frameworks in that title are not asserted as results of the original source extract.
The source catalogue labels the historical question solved. Zhong and Boumal's primary abstract states global convergence of their generalized power method, with linear rate, in a noise regime $\sigma=O(\sqrt{n/\log n})$. Their algorithm and initialization hypotheses must be used; this does not assert convergence of every unshifted iteration from every starting point. The original question and catalogue label are retained separately from the cited result.
\paragraph{Statement recorded in the supplied source.}
Does the projected power method converge (with high probability) to the optimal solution of the angular synchronization problem with (small enough) gaussian noise? \sref{statement}{2}
\paragraph{Source status.}
The supplied catalogue labels this source record solved; that metadata was checked against the supplied snapshot on 2026-10-11. The label is preserved for provenance and is not an independent verification of present open status or mathematical correctness. This is documentation of a source-labelled solved record, not a newly proposed unsolved problem.
\subsection{Short English statement}
Does the specified projected power iteration find a globally optimal phase-synchronization solution when Gaussian noise is sufficiently small?
\subsection{Sources}
\begin{itemize}
\item[\hypertarget{source:statement:1}{S1}] Ulam AI and the UnsolvedMath Contributors, UnsolvedMath: A Curated Collection of Open Mathematics Problems, catalogue entry. Catalogue curation and metadata: CC BY 4.0. The source extract is retained; definitions, English summaries and qualifications are editorial additions.
\url{https://huggingface.co/datasets/ulamai/UnsolvedMath}.
\item[\hypertarget{source:statement:2}{S2}] AMR Open Problem Lists, source linked by the catalogue.
\url{https://afonsobandeira.wordpress.com/2015/12/09/42l10psynchronization/}.
\item[\hypertarget{source:statement:3}{S3}] Yiqiao Zhong and Nicolas Boumal, Near-optimal bounds for phase synchronization, arXiv:1703.06605; SIAM Journal on Optimization, DOI 10.1137/17M1122025. Abstract inspected for the scoped convergence claim.
\url{https://arxiv.org/abs/1703.06605}.
\end{itemize}
\paragraph{Preparation and review.}
The original source and UnsolvedMath contributors retain attribution for the source extract. Codex (GPT-6; exact serving revision not exposed) assisted with definitions, English summaries, scope notes and formatting. No new human mathematical authorship or independently verified solution is asserted. The supplied catalogue credits: Afonso S. Bandeira. Catalogue identity checked against the complete supplied registry and current repository. Distinct upstream records, including repeated questions and broad or damaged extracts, are retained for complete coverage of the source CS category. A repeated record is not a new independent conjecture. Complete mathematical-equivalence and current-literature review has not been performed. Preparation and final review follow contribution prompts C and D, with source gaps disclosed. Structural and rendering checks do not establish a new theorem or open status.
\label{end:statement}
\end{document}
